%% ════════════════════════════════════════════════════════════════════════════
\section{Related Work}
\label{sec:related}

\subsection{Machine Learning for Air Quality Modeling in India}
The application of predictive algorithms to air pollution across Indian metropolitan centres has expanded rapidly. Urban air quality management in India has been reviewed extensively \cite{gulia2015}. Early forecasting studies relied primarily on classical statistical techniques such as autoregressive integrated moving average (ARIMA) models and multiple linear regression. With the public dissemination of continuous ambient air quality monitoring station (CAAQMS) records by the CPCB, researchers transitioned to nonlinear machine learning and deep learning algorithms.

In Gujarat, Maltare and Vahora~\cite{maltare2023} evaluated SARIMA, SVM, and LSTM on Ahmedabad (2015--2021), reporting $R^2 = 0.9989$ for SVM-RBF and $R^2 = 0.951$ for LSTM without specifying a chronological validation split or benchmarking against persistence. At the state level, Advitha et al.~\cite{advitha2024} applied Random Forest and XGBoost to seven Gujarat cities (including Ahmedabad, Ankleshwar, Surat, and Vapi), reporting $R^2 = 98.87\%$ and $98.27\%$ respectively for predicting the calculated AQI from $\text{PM}_{2.5}$, $\text{PM}_{10}$, $\text{SO}_2$, and $\text{NO}_x$ concentrations, without describing a train--test split or a persistence baseline; in their own correlation analysis the calculated AQI is perfectly correlated with $\text{PM}_{10}$ ($r = 1.00$), consistent with an AQI computed deterministically from the model inputs. Surat has also been studied with ARIMA \cite{mahendra2023}. Across broader Indian regions, Sidhu et al.~\cite{sidhu2024} applied CatBoost, XGBoost, and Random Forest across 22 stations in Delhi, Haryana, and Punjab, reporting $R^2 \approx 1.0$ when feeding eight same-day pollutants to predict same-day AQI. Similarly, Chauhan and Sutaria~\cite{chauhan2026} compared SVM, GBM, and XGBoost across six Indian metropolitan cities (including Ahmedabad) using concurrent hourly pollutant inputs, while Mishra et al.~\cite{mishra2025} reported $R^2 = 0.961\text{--}0.9985$ across South Indian cities using a randomly ordered 80:20 split with seven same-day pollutants. Jayapradha et al.~\cite{jayapradha2025} reported $R^2 = 0.948$ for a hybrid LSTM--RF--XGBoost model on multi-city Kaggle records where both same-day pollutants and the derived AQI category (\texttt{AQI\_Bucket}) ranked among top features, and Chakraborty and Shanthi~\cite{chakraborty2026} (a preprint) reported GRU $R^2 = 0.9886$ using AQI-average features whose windows include the current timestep ($i=0$), which risks target leakage if $\text{AQI}_t$ is the target (the target is not explicitly stated). By contrast, Singh et al.~\cite{singh2025} candidly documented severe generalization collapse in pan-India $\text{PM}_{2.5}$ modeling (Random Forest train $R^2 = 0.99$ versus test $R^2 = 0.35$). Meanwhile, Malhotra et al.~\cite{malhotra2024} systematically surveyed AI-based air pollution prediction studies, with an emphasis on India, cataloging validation schemes (e.g., $k$-fold, 70/30) without auditing data leakage.

Close examination shows that several of these studies \cite{sidhu2024}\cite{mishra2025}\cite{jayapradha2025}\cite{chauhan2026}\cite{advitha2024} use concurrent pollutant inputs to predict same-day AQI, and some also use random train-test splits \cite{mishra2025}; for others \cite{maltare2023}\cite{advitha2024}\cite{chakraborty2026}, the validation design is unreported or ambiguous. These are the two leakage mechanisms described in Section~\ref{sec:intro}.

\subsection{Temporal Data Leakage and Validation Protocols}
Data leakage occurs when training data inadvertently incorporates information that would be unavailable at inference time, inducing severe optimistic bias \cite{kaufman2012}. In temporal and spatial machine learning, leakage typically arises through lookahead bias (conditioning on future events) or serial dependency exploitation \cite{bergmeir2012}.

Cerqueira et al.~\cite{cerqueira2020} conducted an extensive comparative benchmark of performance-estimation methods for time series, finding that standard random $k$-fold cross-validation produces systematically over-optimistic estimates for non-stationary data, whereas out-of-sample holdout and prequential (rolling-origin) methods yield estimates that mirror real-world generalization. The severity of this distortion in applied benchmarks is further documented by Albelali and Ahmed~\cite{albelali2025}, who showed that LSTM models trained with cross-validation on sliding windows (built before partitioning) overestimate performance by up to 20.5\% in RMSE, and by Garc{\'i}a Crespi et al.~\cite{garciacrespi2026}, who demonstrated using daily PM$_{10}$ data from Elche, Spain that rolling-origin validation \emph{reverses model rankings} relative to a simple persistence baseline when compared to static hold-out splits. Outside India, Xie and Tummala~\cite{xie2025} benchmarked boosting algorithms across healthcare, environmental, and energy sensor data and found that random $k$-fold cross-validation on Beijing $\text{PM}_{2.5}$ inflated $R^2$ to roughly $0.87$, whereas chronological validation produced negative $R^2$ values. Similarly, Ayus et al.~\cite{ayus2023} reported that XGBoost outperformed deep learning across 10 Chinese cities based on a random 80/20 split with concurrent pollutant inputs, underscoring how pervasive unvalidated splitting remains in atmospheric ML.

\subsection{Persistence Benchmarking and the Forecast Skill Metric}
A fundamental principle of meteorological forecast verification is that any proposed prognostic model must be evaluated against standard climatological and persistence benchmarks \cite{murphy1992}\cite{jolliffe2012}. In strongly autocorrelated processes, the naive Persistence forecaster ($\hat{y}_{t+h} = y_t$) represents a formidable obstacle, particularly at short forecast horizons ($h=1$).

Time-aware evaluation has begun to appear for Indian cities more broadly: Goetzinger et al.~\cite{goetzinger2026} used leakage-safe lag features with a strictly chronological holdout, and Dar et al.~\cite{dar2025} used rolling-origin validation with naive seasonal baselines across ten cities. Most directly relevant to our study, Abidha et al.~\cite{abidha2026} conducted a leakage-aware case study for Chennai (2018--2024), demonstrating that same-day pollutant inputs were near-exact linear rescalings of AQI ($r > 0.9999$); once concurrent inputs were removed, reported $R^2$ collapsed from $1.0$ to $\approx 0.52$, nearly matched by naive Persistence ($R^2 = 0.475$, MAE $= 12.4$). Standard forecast-verification practice recommends reporting skill relative to a reference forecast such as persistence \cite{murphy1992}\cite{jolliffe2012}:
\begin{equation}
  \text{Skill} = 1 - \frac{\text{MAE}_{\text{model}}}{\text{MAE}_{\text{persistence}}}
  \label{eq:skill_def}
\end{equation}
where positive values signify genuine predictive value beyond inertial extrapolation. Extending Abidha et al.~\cite{abidha2026}, our study provides, to our knowledge, the first leak-free, persistence-benchmarked multi-city evaluation for Gujarat (including the industrial corridors of Surat, Vapi, and Ankleshwar, for which we are not aware of prior persistence-referenced evaluation), isolates random-split versus concurrent-input inflation via nested Setups A, B, and C, evaluates multi-day horizon decay ($h=1$ vs.\ $h=3$), and conducts formal Diebold-Mariano hypothesis testing \cite{diebold1995}.

